% GENERATED FILE. Edit canonical lessons, then run npm run generate:books.
% canonical-source-hash: fnv1a64:846752d80fb7bff4
% canonical-lessons: KA-C349-khadga, KA-C349-bana, KA-C349-kavaca, KA-C349-sarathi, KA-C349-madi

\chapter{Swords, Arrows and Armour}
\label{ch:ka-swords-arrows-and-armour}
\hlchaptermodality{349}

\begin{chapteropening}
\textbf{By the end of this chapter:} \emph{I can say a sword, an arrow, armour, a charioteer and to fall in battle.}

Complete the last lesson of chapter 349: I can say a sword, an arrow, armour, a charioteer and to fall in battle.
\end{chapteropening}

\section[\texorpdfstring{\kn{ಖಡ್ಗ} (khaḍga)}{ಖಡ್ಗ (khaḍga)}]{\kn{ಖಡ್ಗ} (khaḍga) — a sword}

\label{lesson:KA-C349-khadga}



\begin{quote}
Before the new one: say the Kannada for a secret, then the Kannada for a sorcerer.
\end{quote}

\subsection*{You'll want to know: \kn{ಖಡ್ಗ}}
\textbf{\kn{ಖಡ್ಗ}} — \emph{khaḍga} — \textquotedblleft{}a sword\textquotedblright{}.

\begin{quote}
\kn{ಆ} \kn{ಯೋಧ} \kn{ಖಡ್ಗ} \kn{ಹಿಡಿದು} \kn{ಮುಂದೆ} \kn{ನಡೆದನು}.
\end{quote}

\begin{quote}
\emph{ā yōdha khaḍga hiḍidu munde naḍedanu.}
\end{quote}

\begin{quote}
That warrior walked ahead, sword in hand.
\end{quote}

Plural: \kn{ಖಡ್ಗಗಳು}, \emph{khaḍgagaḷu}.

The everyday word is \kn{ಕತ್ತಿ}; \kn{ಖಡ್ಗ} belongs to stories and history.

\begin{cousinweb}
\textbf{\kn{ಖಡ್ಗ}} (\emph{khaḍga}) is borrowed from Sanskrit \emph{khaḍga}, \textquotedblleft{}a sword, a scimitar\textquotedblright{} (Monier-Williams).
\end{cousinweb}

\begin{grammarlens}[title={What you've built}]
One more word for telling what happened.
\end{grammarlens}

\subsection*{Your turn}
\begin{itemize}
\raggedright
  \item \emph{Say it:} \emph{khaḍga}
  \item \emph{Say it:} \emph{khaḍga}, once more
  \item \emph{Say it:} \emph{ā yōdha khaḍga hiḍidu munde naḍedanu.} — That warrior walked ahead, sword in hand.
  \item \emph{Recall:} say the Kannada for a secret, then the Kannada for a sorcerer, then say \emph{khaḍga} again
\end{itemize}

\begin{tcolorbox}[breakable,colback=teal!4,colframe=teal!35!black,title={Before you move on}]
What does \kn{ಖಡ್ಗ} mean? (\textquotedblleft{}A sword\textquotedblright{}.) Say it once more.
\end{tcolorbox}
\section[\texorpdfstring{\kn{ಬಾಣ} (bāṇa)}{ಬಾಣ (bāṇa)}]{\kn{ಬಾಣ} (bāṇa) — an arrow}

\label{lesson:KA-C349-bana}



\begin{quote}
Before the new one: say the Kannada for a sorcerer, then the Kannada for a sword.
\end{quote}

\subsection*{You'll want to know: \kn{ಬಾಣ}}
\textbf{\kn{ಬಾಣ}} — \emph{bāṇa} — \textquotedblleft{}an arrow\textquotedblright{}.

\begin{quote}
\kn{ಬೇಟೆಗಾರನ} \kn{ಬಾಣ} \kn{ಮರಕ್ಕೆ} \kn{ತಾಗಿತು}.
\end{quote}

\begin{quote}
\emph{bēṭegārana bāṇa marakke tāgitu.}
\end{quote}

\begin{quote}
The hunter's arrow hit the tree.
\end{quote}

Plural: \kn{ಬಾಣಗಳು}, \emph{bāṇagaḷu}.

A bow is \kn{ಬಿಲ್ಲು}; arrows are shot: \kn{ಬಾಣ} \kn{ಬಿಡು}.

\begin{cousinweb}
\textbf{\kn{ಬಾಣ}} (\emph{bāṇa}) is borrowed from Sanskrit \emph{bāṇa}, \textquotedblleft{}a reed shaft, an arrow\textquotedblright{} (CDIAL 9203).
\end{cousinweb}

\begin{grammarlens}[title={What you've built}]
One more word for telling what happened.
\end{grammarlens}

\subsection*{Your turn}
\begin{itemize}
\raggedright
  \item \emph{Say it:} \emph{bāṇa}
  \item \emph{Say it:} \emph{bāṇa}, once more
  \item \emph{Say it:} \emph{bēṭegārana bāṇa marakke tāgitu.} — The hunter's arrow hit the tree.
  \item \emph{Recall:} say the Kannada for a sorcerer, then the Kannada for a sword, then say \emph{bāṇa} again
\end{itemize}

\begin{tcolorbox}[breakable,colback=teal!4,colframe=teal!35!black,title={Before you move on}]
What does \kn{ಬಾಣ} mean? (\textquotedblleft{}An arrow\textquotedblright{}.) Say it once more.
\end{tcolorbox}
\section[\texorpdfstring{\kn{ಕವಚ} (kavaca)}{ಕವಚ (kavaca)}]{\kn{ಕವಚ} (kavaca) — armour}

\label{lesson:KA-C349-kavaca}



\begin{quote}
Before the new one: say the Kannada for a sword, then the Kannada for an arrow.
\end{quote}

\subsection*{You'll want to know: \kn{ಕವಚ}}
\textbf{\kn{ಕವಚ}} — \emph{kavaca} — \textquotedblleft{}armour\textquotedblright{}.

\begin{quote}
\kn{ಸೈನಿಕರು} \kn{ಕಬ್ಬಿಣದ} \kn{ಕವಚ} \kn{ತೊಟ್ಟಿದ್ದರು}.
\end{quote}

\begin{quote}
\emph{sainikaru kabbiṇada kavaca toṭṭiddaru.}
\end{quote}

\begin{quote}
The soldiers wore iron armour.
\end{quote}

Also any protective cover: a phone case is a \kn{ಕವಚ} too.

\begin{cousinweb}
\textbf{\kn{ಕವಚ}} (\emph{kavaca}) is borrowed from Sanskrit \emph{kavaca}, \textquotedblleft{}a coat of mail, a jacket\textquotedblright{} (CDIAL 2957).
\end{cousinweb}

\begin{grammarlens}[title={What you've built}]
One more word for telling what happened.
\end{grammarlens}

\subsection*{Your turn}
\begin{itemize}
\raggedright
  \item \emph{Say it:} \emph{kavaca}
  \item \emph{Say it:} \emph{kavaca}, once more
  \item \emph{Say it:} \emph{sainikaru kabbiṇada kavaca toṭṭiddaru.} — The soldiers wore iron armour.
  \item \emph{Recall:} say the Kannada for a sword, then the Kannada for an arrow, then say \emph{kavaca} again
\end{itemize}

\begin{tcolorbox}[breakable,colback=teal!4,colframe=teal!35!black,title={Before you move on}]
What does \kn{ಕವಚ} mean? (\textquotedblleft{}Armour\textquotedblright{}.) Say it once more.
\end{tcolorbox}
\section[\texorpdfstring{\kn{ಸಾರಥಿ} (sārathi)}{ಸಾರಥಿ (sārathi)}]{\kn{ಸಾರಥಿ} (sārathi) — a charioteer}

\label{lesson:KA-C349-sarathi}



\begin{quote}
Before the new one: say the Kannada for an arrow, then the Kannada for armour.
\end{quote}

\subsection*{You'll want to know: \kn{ಸಾರಥಿ}}
\textbf{\kn{ಸಾರಥಿ}} — \emph{sārathi} — \textquotedblleft{}a charioteer\textquotedblright{}.

\begin{quote}
\kn{ಸಾರಥಿ} \kn{ರಥವನ್ನು} \kn{ವೇಗವಾಗಿ} \kn{ಓಡಿಸಿದನು}.
\end{quote}

\begin{quote}
\emph{sārathi rathavannu vēgavāgi ōḍisidanu.}
\end{quote}

\begin{quote}
The charioteer drove the chariot fast.
\end{quote}

Plural: \kn{ಸಾರಥಿಗಳು}, \emph{sārathigaḷu}.

Today also a leader who steers a group: \kn{ತಂಡದ} \kn{ಸಾರಥಿ}.

\begin{cousinweb}
\textbf{\kn{ಸಾರಥಿ}} (\emph{sārathi}) is borrowed from Sanskrit \emph{sārathi}, \textquotedblleft{}a charioteer, the driver of a car\textquotedblright{} (Monier-Williams), from \emph{sa-ratha}, \textquotedblleft{}with the chariot\textquotedblright{}.
\end{cousinweb}

\begin{grammarlens}[title={What you've built}]
One more word for telling what happened.
\end{grammarlens}

\subsection*{Your turn}
\begin{itemize}
\raggedright
  \item \emph{Say it:} \emph{sārathi}
  \item \emph{Say it:} \emph{sārathi}, once more
  \item \emph{Say it:} \emph{sārathi rathavannu vēgavāgi ōḍisidanu.} — The charioteer drove the chariot fast.
  \item \emph{Recall:} say the Kannada for an arrow, then the Kannada for armour, then say \emph{sārathi} again
\end{itemize}

\begin{tcolorbox}[breakable,colback=teal!4,colframe=teal!35!black,title={Before you move on}]
What does \kn{ಸಾರಥಿ} mean? (\textquotedblleft{}A charioteer\textquotedblright{}.) Say it once more.
\end{tcolorbox}
\section[\texorpdfstring{\kn{ಮಡಿ} (maḍi)}{ಮಡಿ (maḍi)}]{\kn{ಮಡಿ} (maḍi) — to fall in battle; to die}

\label{lesson:KA-C349-madi}



\begin{quote}
Before the new one: say the Kannada for armour, then the Kannada for a charioteer.
\end{quote}

\subsection*{You'll want to know: \kn{ಮಡಿ}}
\textbf{\kn{ಮಡಿ}} — \emph{maḍi} — \textquotedblleft{}to fall in battle; to die\textquotedblright{}.

\begin{quote}
\kn{ಆ} \kn{ಯುದ್ಧದಲ್ಲಿ} \kn{ಸಾವಿರಾರು} \kn{ಸೈನಿಕರು} \kn{ಮಡಿದರು}.
\end{quote}

\begin{quote}
\emph{ā yuddhadalli sāvirāru sainikaru maḍidaru.}
\end{quote}

\begin{quote}
Thousands of soldiers fell in that war.
\end{quote}

Past (I …): \kn{ಮಡಿದೆ}, \emph{maḍide}.

A respectful verb for dying, used of soldiers and heroes. A different \kn{ಮಡಿ} means ritual cleanness.

\begin{cousinweb}
\textbf{\kn{ಮಡಿ}} (\emph{maḍi}) is native Dravidian: DEDR 4653 lists Kannada \emph{maḍi}, \textquotedblleft{}to die\textquotedblright{}, and \emph{maḍuhu}, \textquotedblleft{}to kill\textquotedblright{}.
\end{cousinweb}

\begin{grammarlens}[title={What you've built}]
One more word for telling what happened.
\end{grammarlens}

\subsection*{Your turn}
\begin{itemize}
\raggedright
  \item \emph{Say it:} \emph{maḍi}
  \item \emph{Say it:} \emph{maḍi}, once more
  \item \emph{Say it:} \emph{ā yuddhadalli sāvirāru sainikaru maḍidaru.} — Thousands of soldiers fell in that war.
  \item \emph{Recall:} say the Kannada for armour, then the Kannada for a charioteer, then say \emph{maḍi} again
\end{itemize}

\begin{tcolorbox}[breakable,colback=teal!4,colframe=teal!35!black,title={Before you move on}]
What does \kn{ಮಡಿ} mean? (\textquotedblleft{}To fall in battle; to die\textquotedblright{}.) Say it once more.
\end{tcolorbox}
